\documentclass{article}
\usepackage[utf8]{inputenc}
\usepackage[T1]{fontenc}
\usepackage{amsmath, amssymb, amsthm}
\usepackage{geometry}
\usepackage{hyperref}
\usepackage{bookmark}
\usepackage{mdframed}
\usepackage{xcolor}
\usepackage{amsmath}


% \renewcommand{\thesubsection}{\Roman{subsection}}
% \newcommand{\exo}[2]{\textbf{\Roman{#1} #2} }

% \geometry{left=2.5cm, right=2.5cm, top=2.5cm, bottom=2.5cm}


\setlength{\parindent}{0pt}
\setlength{\parskip}{0.8em}

\title{Intégration}
\author{Axel}
\date{\today}




\newcounter{exonb}
\stepcounter{exonb} 
\newcommand{\exo}[2]{
    \vspace{5mm}
    \textbf{#1}
    \textbf{\arabic{exonb}}
    \hrulefill
    \textbf{#2}
    \stepcounter{exonb} 
}


\newcounter{prblmnb}
\stepcounter{prblmnb} 
\newcommand{\problem}[1]{
    \vspace{5mm}
    \textbf{Problème}
    \textbf{\Roman{prblmnb}}
    \hrulefill
    \textbf{#1}
    \stepcounter{prblmnb} 
}


\newcommand{\smalltext}[1]{
    \scriptsize
    \textit{#1}
    \normalsize
}





\newenvironment{correction}
    {\begin{mdframed}[backgroundcolor=gray!20, roundcorner=10pt, linecolor=gray!10,
        leftmargin=0, rightmargin=0, innerleftmargin=5, innerrightmargin=5,
        innertopmargin=0, innerbottommargin=10]}
    {\end{mdframed}}
\begin{document}

\newcommand{\correct}[1]{
    \begin{correction}
        #1
    \end{correction}
}


\newcommand{\exocor}[3]{
    \exo{Exercice}{#3}

    \begin{equation*}
        #1
    \end{equation*}

    \correct{
        \begin{gather*}
            #2
        \end{gather*}
    }
}

\maketitle

les exercices sont classé par difficultée
les étoiles(\star) indiquent la difficultée de l'exercice
Les exercices sont réparties par leurs différentes methodes de résolution, 
puis des problèmes plus complexes sont proposés à la fin. demandant une combinaison de méthodes de résolution.

% les symboles suivantes sont utiliser pour proposer une Methode de resolution
% \begin{itemize}
%     \item[--] \odot  \: Forme évidente
%     \item[--] \oslash   \: méthode de décomposition en éléments simples
%     \item[--] \circledcirc   \: changement de variable
%     \item[--] \otimes  \: intégration par partie
% \end{itemize}

Nous donnerons une solution pour chaque exercice.


\section*{Exercices}
\subsection*{Formes usuels}

\exo{Exercice}{}
\begin{equation*}
    \int \frac{1}{x^2} \, dx
\end{equation*}

\correct{
    \begin{gather*}
        -\frac{1}{x}
    \end{gather*}
}

\exo{Exercice}{}
\begin{equation*}
    \int \sqrt{x} \, dx
\end{equation*}

\correct{
    \begin{gather*}
        \frac{2}{3} x^{3/2}
    \end{gather*}
}

\exo{Exercice}{}
\begin{equation*}
    \int \frac{1}{x+1} \, dx
\end{equation*}


\correct{
    \begin{gather*}
        \ln |x + 1|
    \end{gather*}
}

\exo{Exercice}{}
\begin{equation*}
    \int \frac{1}{x^2+4} \, dx
\end{equation*}

\correct{
    \begin{gather*}
        \frac{1}{2} \arctan \left( \frac{x}{2} \right)
    \end{gather*}
}

\exo{Exercice}{}
\begin{equation*}
    \int 2x e^{x^2} \, dx
\end{equation*}

\correct{
    \begin{gather*}
        e^{x^2}
    \end{gather*}
}

\exo{Exercice}{}
\begin{equation*}
    \int \frac{1}{\sqrt{2-x^2}} \, dx
\end{equation*}

\correct{
    \begin{gather*}
        \arcsin \left( \frac{x}{\sqrt{2}} \right)
    \end{gather*}
}



\exo{Exercice}{}
\begin{equation*}
    \int 2x\cos(x^2) \, dx
\end{equation*}

\correct{
    \begin{gather*}
        \sin(x^2)
    \end{gather*}
}

\exo{Exercice}{\star}
\begin{equation*}
    \int \sin(x)\cos(\cos(x)) \, dx
\end{equation*}

\correct{
    \begin{gather*}
        -\sin(\cos(x))
    \end{gather*}
}


\subsection*{Fractions simples}


\exo{Exercice}{\star}

\begin{equation*}
    \int \frac{1}{x^2 - 1} \, dx
\end{equation*}

\correct{
    \begin{gather*}
        \ln \left| \frac{\sqrt{x - 1}}{\sqrt{x + 1}} \right|
    \end{gather*}
}

\exo{Exercice}{}

\begin{equation*}
    \int \frac{x}{x^2 + 1} \, dx
\end{equation*}

\correct{
    \begin{gather*}
        \frac{1}{2} \ln |x^2 + 1|
    \end{gather*}
}


\exo{Exercice}{}

\begin{equation*}
    \int \frac{1}{x^2 + x} \, dx
\end{equation*}

\correct{
    \begin{gather*}
        \ln |x| - \ln |x + 1|
    \end{gather*}
}

\subsection*{Logarithmes}

\exo{Exercice}{\star}

\begin{equation*}
    \int xln(x) \, dx
\end{equation*}

\correct{
    \begin{gather*}
        \frac{x^2}{2} \ln |x| - \frac{x^2}{4}
    \end{gather*}
}

\exo{Exercice}{\star\star}

\begin{equation*}
    \int x^2 ln(x) \, dx
\end{equation*}

\correct{
    \begin{gather*}
        \frac{1}{9} x^3 (3 \ln x - 1)
    \end{gather*}
}

\exo{Exercice}{\star\star}

\begin{equation*}
    \int ln(x)^2 \, dx
\end{equation*}

\correct{
    \begin{gather*}
        x (ln(x)^2 - 2ln(x) + 2)
    \end{gather*}
}


\subsection*{Polynomes exponentielles}

\exo{Exercice}{}

\begin{equation*}
    \int xe^{x} \, dx
\end{equation*}

\correct{
    \begin{gather*}
        (x - 1) e^x
    \end{gather*}
}

\exo{Exercice}{\star}

\begin{equation*}
    \int x^2 e^x \, dx
\end{equation*}

\correct{
    \begin{gather*}
        (x^2 - 2x + 2) e^x
    \end{gather*}
}

\exo{Exercice}{\star}

\begin{equation*}
    \int x^3 e^x \, dx
\end{equation*}

\correct{
    \begin{gather*}
        (x^3 - 3x^2 + 6x - 6) e^x
    \end{gather*}
}

\subsection*{Trigo exponentielles}

\exocor{
    \int e^x \sin(x) \, dx
}{
    \frac{1}{2} e^x (\sin(x) - \cos(x))
}{\star\star}

\exocor{
    \int sin(x^2) e^{\frac{x^2}{2}} dx
}{
    \frac{1}{5} e^{x^2/2} (sin(x^2) - 2 cos(x^2))
}{\star\star\star}


\problem{\star\star\star}

Soit $\omega, \mu \in \mathbb{R}^*$

$(1)$ Donner l'expression de l'intégrale suivante:

\begin{equation*}
    \int \sin(\omega x + \mu ) \, dx
\end{equation*}

$(2)$ En déduire l'expression de l'intégrale suivante:

\begin{equation*}
    \int e^x \sin(\omega x + \mu ) \, dx
\end{equation*}

$(3)$ Soit $\alpha \in \mathbb{R}^*$ Donner l'expression de l'intégrale suivante:

\begin{equation*}
    \int e^{\alpha x} \sin(\omega x + \mu ) \, dx
\end{equation*}

\correct{
    \begin{gather*}
        -\frac{1}{\omega} \cos(\omega x + \mu) \\
        \frac{e^x}{\omega^2 + 1} (\sin(\omega x + \mu) - \omega \cos(\omega x + \mu)) \\
        \dots
    \end{gather*}
}

\subsection*{Fraction rationnelle}

\exocor{
    \int \frac{1}{x^2-x-1} \, dx
}{
    \frac{\ln\left(-2x+\sqrt{5}+1\right)-\ln\left(2x+\sqrt{5}-1\right)}{\sqrt{5}}
}{\star\star\star}

\exocor{
    \int \frac{1}{x^2+x+1} \, dx
}{
    \frac{2\arctan\left(\frac{2x+1}{\sqrt{3}}\right)}{\sqrt{3}}
}{\star\star}

\exocor{
    \int \frac{1}{x^3+1} \, dx
}{
    \frac{1}{6}\left[2\log(x+1)-\ln\left(x^{2}-x+1\right)+2\sqrt{3}\arctan\left(\frac{2x-1}{\sqrt{3}}\right)\right]
}{\star\star\star\star}

\subsection*{Introduction aux règles de Bioche}

\exocor{
    \int \sin^{3}(x) \, dx
}{
    -\cos(x) + \frac{\cos^3(x)}{3} 
}{\star}

\exocor{
    \int \cos^{3}(x) \, dx
}{
    \sin(x) - \frac{\sin^3(x)}{3}
}{\star}

\exocor{
    \int \tan^{3}(x) \, dx
}{
    \frac{1}{2\cos^2(x)}+\ln(\cos(x))
}{\star\star}

\exocor{
    \int \frac{\sin^3(x)}{1+\cos^2x} \, dx
}{
    \frac{1}{2\cos^2(x)}+\ln(\cos(x))
}{\star\star}

\exocor{
    \int \frac{\cos(t)}{1+\sin(t)} \, dx
}{
    \ln{\left|1+\sin(t)\right|}
}{\star\star}



\end{document}


